\section{Kritéria výběru modelu}
\label{sec:selection-criteria}

Výběr vhodného jazykového modelu pro integraci do systému Kelvin není triviální rozhodnutí.
Na trhu existují desítky modelů s různými vlastnostmi, způsoby nasazení a cenovými modely.
Aby bylo možné provést strukturované srovnání, je nutné nejprve definovat kritéria, podle nichž budou modely posuzovány.
Tato sekce popisuje čtyři klíčová kritéria relevantní pro konkrétní použití v kontextu analýzy studentských programátorských úloh.

\subsection{Kvalita výstupů}
\label{subsec:criteria-quality}

Nejdůležitějším kritériem je schopnost modelu generovat výstupy, které jsou pro učitele i studenty skutečně užitečné.
Kvalita výstupů v kontextu analýzy kódu zahrnuje několik dílčích vlastností:

\begin{itemize}
    \item \textbf{Věcná správnost} -- model musí identifikovat reálné problémy v kódu a nepřidávat falešné výstrahy (\emph{false positives}).
    Falešný pozitivní nález (tedy upozornění na problém, který neexistuje) je pro studenta matoucí a pro učitele zbytečná zátěž.

    \item \textbf{Relevance k zadání} -- model by měl posuzovat řešení s ohledem na konkrétní zadání úlohy, nejen na obecné programátorské principy.
    Například to, zda je zvolený algoritmus přiměřený požadovaným omezením (velikost vstupu, požadavky na rychlost), není z kódu samotného zjistitelné bez kontextu zadání.

    \item \textbf{Srozumitelnost a didaktická hodnota} -- výstupy musí být psány jazykem, kterému studenti rozumí.
    Technické termíny jsou v pořádku, pokud jsou používány přesně, ale výstup nesmí být prázdnou formálností.
    Model by měl vysvětlit \emph{proč} je daná část kódu problematická a navrhnout konkrétní způsob opravy.

    \item \textbf{Konzistentnost} -- model by měl pro podobné vstupy generovat podobné výstupy.
    Absolutní reprodukovatelnost není u pravděpodobnostních systémů realistická, ale výstupy by se neměly mezi dvěma identickými odevzdáními zásadně lišit.
\end{itemize}

Pro objektivní porovnání kvality modelů existují standardizované benchmarky\footnote{Benchmark je soubor testovacích úloh sloužící k měření výkonu systému za kontrolovaných podmínek.}
zaměřené na schopnost porozumění a generování kódu.
Nejznámější z nich jsou:

\begin{itemize}
    \item \textbf{HumanEval} \cite{humaneval2021} -- sada 164 programátorských úloh v jazyce Python s automatickým ověřením správnosti generovaného kódu.
    Výsledek (pass@1) udává, jak často model na první pokus vygeneruje správné řešení.

    \item \textbf{MBPP} (Mostly Basic Python Problems) \cite{mbpp2021} -- podobný benchmark s důrazem na jednoduché úlohy vhodné pro začátečníky.

    \item \textbf{LiveCodeBench} \cite{livecodebench2024} -- novější benchmark zahrnující úlohy z programátorských soutěží, které průběžně přibývají,
    čímž snižuje riziko kontaminace trénovacích dat\footnote{Kontaminace trénovacích dat nastane, když jsou testovací úlohy obsaženy v datech, na nichž byl model trénován. V takovém případě model úlohu nemusí skutečně \uv{vyřešit} -- mohl ji memorovat.}.
\end{itemize}

Je důležité mít na paměti, že benchmarky měří především schopnost \emph{generovat} kód, nikoli schopnost \emph{analyzovat} a \emph{komentovat} existující kód.
Pro systém Kelvin je cílem druhá jmenovaná schopnost.
Benchmarky slouží jako orientační ukazatel kvality modelu, ale nelze z nich přímo vyvodit závěry o kvalitě generovaných code review komentářů.
Proto jsou výsledky benchmarků v této kapitole doplněny o kvalitativní hodnocení na základě manuálního testování vybraných modelů na vzorku historických studentských odevzdání.

\subsection{Dostupnost modelu}
\label{subsec:criteria-availability}

Dostupností se v tomto kontextu rozumí několik vzájemně propojených otázek:

\begin{itemize}
    \item \textbf{Způsob přístupu} -- je model dostupný jako cloudové API, nebo jako open-source model ke stažení a lokálnímu nasazení?
    Cloudové API nabízí snadnou integraci bez hardwarových nároků, ale zavádí závislost na externím poskytovateli.
    Open-source modely umožňují lokální nasazení a plnou kontrolu nad zpracováním dat.

    \item \textbf{Licenční podmínky} -- pro akademické a potenciálně komerční použití v rámci instituce je nutné ověřit, zda licence modelu takové použití umožňuje.
    Mnoho modelů je dostupných pod permisivními licencemi (Apache~2.0, MIT), ale některé mají omezení pro komerční použití nebo vyžadují uzavření samostatné licenční smlouvy.

    \item \textbf{Geografická dostupnost a datová centra} -- pro českou akademickou instituci je relevantní, zda cloudová API zpracovávají data v rámci Evropského hospodářského prostoru (EHP).
    Přenos osobních údajů studentů mimo EHP podléhá přísným podmínkám dle nařízení GDPR.

    \item \textbf{Stabilita a podpora} -- cloudové API musí být spolehlivé a výrobce musí poskytovat průběžnou podporu.
    Pro akademické použití je důležitá i předvídatelnost cenového modelu a garancí dostupnosti (SLA).
\end{itemize}

\subsection{Náročnost provozu}
\label{subsec:criteria-operational-cost}

Náročnost provozu určuje, jaké zdroje jsou potřeba k dlouhodobému provozování modelu v produkčním prostředí.
Zahrnuje jak finanční náklady, tak nároky na infrastrukturu a údržbu.

U cloudových API jsou náklady typicky určeny cenovými modely \emph{za token}\footnote{Token je základní jednotka textu, s níž modely pracují -- přibližně odpovídá čtyřem znakům nebo třičtvrtě slova v angličtině. Viz sekce~\ref{subsec:llm-tokens-parameters-context}.}.
Pro typickou analýzu jednoho studentského odevzdání, zahrnující systémový prompt, zadání a zdrojový kód, je potřeba přibližně 2\,000--8\,000 vstupních tokenů a 500--2\,000 výstupních tokenů.
Celkové náklady na jedno odevzdání tak závisí na zvoleném modelu a jeho cenové politice.

Pro on-premise\footnote{On-premise (zkráceně on-prem) označuje řešení provozované na vlastní infrastruktuře instituce, nikoliv u externího poskytovatele.} nasazení jsou jednorázové náklady na pořízení GPU\footnote{GPU (Graphics Processing Unit) -- grafický procesor původně navržený pro zobrazování grafiky, ale díky svým masivně paralelním výpočetním schopnostem velmi vhodný i pro trénování a inferenci neuronových sítí.}
serveru výrazně vyšší, avšak průběžné náklady (elektřina, chlazení, správa) jsou nižší než opakující se poplatky za API přístupy.

\subsection{Schopnost pojmout celé studentské řešení}
\label{subsec:criteria-context-window}

Jak bylo popsáno v sekci~\ref{subsec:llm-token-limit}, každý model je omezen maximální délkou vstupu -- takzvaným \emph{kontextovým oknem}.
Pro analýzu studentského řešení musí kontext pojmout:

\begin{enumerate}
    \item \textbf{Systémový prompt} -- instrukce definující roli modelu, požadovaný formát výstupu a obecné pokyny pro analýzu.
    Typicky 200--800 tokenů.

    \item \textbf{Zadání úlohy} -- popis úlohy, požadavky, vstupy a výstupy.
    Typicky 100--500 tokenů, v případě obsáhlého zadání i více.

    \item \textbf{Zdrojový kód studenta} -- samotné odevzdané řešení.
    Pro typické úlohy z výuky programování v jazyce C++ se rozsah pohybuje od 50 do 500 řádků kódu.
    Odpovídá přibližně 500 až 5\,000 tokenům (zdrojový kód je z hlediska tokenizace efektivnější než přirozený jazyk, protože identifikátory a klíčová slova jsou kratší).
\end{enumerate}

Pro standardní studentské odevzdání v systému Kelvin je tedy dostačující kontextové okno přibližně \textbf{8\,000 tokenů}.
Aby byl systém schopen zvládnout i rozsáhlejší úlohy (větší projekty, vícesouborová řešení) a aby byl respektován princip dostatečné rezervy, je jako \textbf{minimální požadavek} stanoveno \textbf{16\,000 tokenů}.
Moderní modely tuto hranici zpravidla bez problémů splňují -- většina aktuálně dostupných modelů nabízí kontextové okno 32\,000 až 128\,000 tokenů.

\begin{table}[H]
    \centering
    \resizebox{\textwidth}{!}{%
        \begin{tabular}{lp{9cm}l}
            \toprule
            \textbf{Kritérium}              & \textbf{Popis}                                                                 & \textbf{Minimální požadavek} \\
            \midrule
            Kvalita výstupů                 & Věcně správné, srozumitelné, relevantní komentáře                              & Vyšší než průměr \\
            Dostupnost                      & Přístupný přes API nebo ke stažení; vhodná licence pro akademické použití      & Open-source nebo API s DPA \\
            Náročnost provozu               & Přijatelné náklady na token nebo na hardware                                   & Viz tabulka~\ref{tab:deployment-comparison} \\
            Kontextové okno                 & Schopnost zpracovat celé odevzdání včetně promptu                              & $\geq$ 16\,000 tokenů \\
            \bottomrule
        \end{tabular}
    }
    \caption{Přehled kritérií výběru modelu a jejich minimálních požadavků}
    \label{tab:selection-criteria}
\end{table}

\endinput
